\chapter{Generative Models for Dynamical Systems}
\label{ch:generative}

\epigraph{\itshape The future of simulation is sampling.}{(Folk wisdom,
SciML circles, c.\ 2024)}

\section{Why generative?}
Many physical systems are inherently stochastic (turbulence,
multiphase flows, ensemble forecasts, climate). Even deterministic
PDEs benefit from probabilistic surrogates that quantify uncertainty
and recover multimodal distributions. Generative models provide
\emph{distributions over solutions} rather than point predictions.

\section{Diffusion models}
\subsection{Score-based formulation}
Define a forward SDE that perturbs data into noise:
\begin{equation}
\dd \bm{X}_t = \bm{f}(\bm{X}_t,t)\dd t + g(t)\dd \bm{W}_t,
\qquad t\in[0,T].
\end{equation}
The reverse SDE generating samples from $p_T$ back to $p_0$ is
\begin{equation}
\dd \bm{X}_t = \big[\bm{f}(\bm{X}_t,t) - g(t)^2\nabla\log p_t(\bm{X}_t)\big]\dd t + g(t)\dd \bar{\bm{W}}_t.
\end{equation}
The score $\nabla\log p_t$ is learned by score
matching~\cite{song2021scoresde}. Equivalent ODE formulation
(``probability flow ODE'') gives deterministic sampling.

\subsection{Diffusion for PDEs}
Diffusion models are increasingly used as PDE surrogates and
super-resolvers:
\begin{itemize}
\item \textbf{PDE-Refiner}~\cite{lippe2023pderefiner}: iterative noise
denoising at multiple scales fixes spectral defects of single-shot
operators.
\item \textbf{DYffusion}~\cite{cachay2023dyffusion}: temporal diffusion
models for stochastic forecasting.
\item \textbf{SimDiffPDE}~\cite{simdiffpde2025}: image-to-image
transformer diffusion as a unified PDE solver.
\item \textbf{Diffusion priors with physics-guided
inference}~\cite{diffphys2026}: train a generic diffusion prior; condition
at inference on PDE residual via guidance.
\item \textbf{GenCast}~\cite{price2024gencast}: diffusion-based ensemble
weather forecasting outperforming top operational systems.
\end{itemize}

\subsection{Conditional / guided diffusion}
\begin{itemize}
\item \textbf{Classifier-free guidance}: train conditional and
unconditional networks jointly.
\item \textbf{Plug-and-play physics guidance}: at sampling, add a
physics-based correction term $-\lambda \nabla \|\Nop[\hat u]\|^2$ to
the score.
\item \textbf{Posterior sampling for inverse problems} (DPS, $\Pi$GDM):
combine diffusion priors with likelihoods $p(\bm{y}|\bm{x})$ for
data assimilation, super-resolution, sparse-view reconstruction.
\end{itemize}

\section{Flow matching and stochastic interpolants}
\paragraph{Flow matching}~\cite{lipman2023flowmatching}.
Choose an interpolation $\bm{x}_t = \alpha_t\bm{x}_1 + \sigma_t\bm{x}_0$
between data $\bm{x}_1$ and noise $\bm{x}_0$, with target velocity
$\bm{u}_t(\bm{x})=\E[\dot{\bm{x}}_t|\bm{x}_t=\bm{x}]$. Train
$\bm{v}_\theta(\bm{x},t)$ by regression on conditional velocities ---
no SDE simulation, no score estimation.
\paragraph{Stochastic interpolants}~\cite{albergo2023interpolants}.
Generalise flow matching with arbitrary interpolation paths and noise
schedules.
\paragraph{Physics-Based Flow Matching (PBFM)~\cite{pbfm2025}.}
Augment the flow-matching loss with PDE residuals and algebraic
constraints, so the learned velocity field generates physically
consistent samples.

\section{Generative ensembles for weather and climate}
A 2024--2025 milestone: ML-based ensemble systems (GenCast, NeuralGCM
ensembles) match or surpass operational ensembles in skill while
running orders of magnitude faster on accelerators. Implications for
SciML:
\begin{itemize}
\item Foundation operators (\cref{ch:emerging}) train deterministic
surrogates; diffusion sits on top to provide ensemble spread.
\item Calibration and reliability metrics (CRPS, rank histograms)
become as important as RMSE.
\end{itemize}

\section{GANs and energy-based models in SciML}
While diffusion has overtaken GANs, mention:
\begin{itemize}
\item \textbf{Physics-informed GANs}: residual losses on the generator
output for super-resolution of turbulence, downscaling.
\item \textbf{Energy-based models / score networks}: useful when the
unnormalised density is the natural object (statistical mechanics,
molecular dynamics free energy).
\end{itemize}

\section{Practical considerations}
\begin{recipe}{Training a diffusion-based PDE surrogate}
\begin{enumerate}
\item Choose noise schedule (variance-preserving, variance-exploding,
EDM/Karras parameterisation).
\item Architectural backbone: U-Net for grids, transformers / FNO for
heterogeneous data.
\item Condition on PDE coefficients, BCs, ICs, time index.
\item Loss: $\bm{v}$-prediction or $\epsilon$-prediction; physics
residual penalty optional.
\item Sampling: deterministic ODE (DDIM, EDM) for speed, SDE for
diversity; few-step distillation for inference.
\end{enumerate}
\end{recipe}

\begin{pitfall}{Generative samples that are not physical}
A diffusion model that scores well on pixel-MSE may still produce
samples violating mass conservation, divergence-free constraints, or
boundary conditions. Add hard projections onto the constraint manifold
(divergence-free projection for incompressible flows) or physics-guided
sampling.
\end{pitfall}
